\documentclass[12pt]{article}
\usepackage[utf8]{inputenc}
\usepackage[T1]{fontenc}
\usepackage{hyperref}
\title{M\'aku Language Resources for K-12 Curriculum Development}
\author{Compiled by Claude Sonnet 5.5 (claude-sonnet-5-5) for Andr\'as Kornai}
\date{}
\begin{document}
\maketitle

\section{Introduction}
M\'aku (also called Jukude; ISO 639-3 xak, Glottocode maku1246) is an extinct language, probably an isolate. It was spoken along the upper Uraricoera and lower Auari rivers in Roraima, near the Brazil--Venezuela border. The last speaker, Sinfr\^onio Magalh\~aes, died in 2000, and no one identifies as Jukude today. Aryon Rodrigues, Ernesto Migliazza and Iraguacema Lima Maciel worked on the language. Do not confuse it with the many unrelated languages called ``Maku'' as an exonym.

\section{Primary Educational Resources}
None found. No school materials or revitalization programme were found.

\section{Linguistic and Reference Materials}
\begin{itemize}
\item Chris Rogers, \emph{M\'aku: A Comprehensive Grammar} (2020). It includes a M\'aku--English--Portuguese glossary and an English--M\'aku word list. The glossary contents were not seen online. Catalogue listing: \url{https://www.sack.de/rogers-mku/9780367522193}
\item M\'aku language of Auari article: \url{https://en.wikipedia.org/wiki/M%C3%A1ku_language_of_Auari}
\item Disambiguation of the exonym: \url{https://en.wikipedia.org/wiki/Maku_(exonym)}
\end{itemize}

\section{Religious and Community Resources}
None found.

\section{Digital and Audio Resources}
None found.

\section{Curriculum Development Context}
The main documentary basis is Rogers's grammar and glossary. A heritage curriculum would have to be built from that book, with its copyright and availability checked, and with community consultation. I found no descendant community that claims the language.
\end{document}